\documentclass[11pt]{article}

\usepackage[T1]{fontenc}
\usepackage{lmodern}
\usepackage[margin=1in]{geometry}
\usepackage{booktabs}
\usepackage{amsmath}
\usepackage{hyperref}
\usepackage{xcolor}
\usepackage{listings}
\usepackage{enumitem}

\lstset{
    basicstyle=\ttfamily\small,
    breaklines=true,
    columns=fullflexible,
    keywordstyle=\color{blue!70!black}\bfseries,
    commentstyle=\color{gray},
    stringstyle=\color{orange!80!black},
    language=Python,
    showstringspaces=false,
    frame=single,
    framesep=4pt,
    xleftmargin=2pt,
    xrightmargin=2pt,
}

\newcommand{\code}[1]{\texttt{#1}}

\title{Supplementary Material:\\Dataset Preprocessing and Concept Bottleneck Model Training}
\author{}
\date{}

\begin{document}
\maketitle

\section{Overview}

This document describes, for both experiments in the study (a phishing-email
detection task and a bird-species classification task on CUB-200-2011),
every implementation step from raw input data to the final task prediction:
dataset preprocessing, encoding, train/validation/test splitting, and the
training of the two-stage Concept Bottleneck Model (CBM) itself. Both
experiments follow the same two-stage CBM design:
\begin{enumerate}
    \item \textbf{Concept extractor} $g_c$: encoder embedding $\rightarrow$
    predicted concept activations.
    \item \textbf{Label predictor} $g_y$: concept activations $\rightarrow$
    predicted task label.
\end{enumerate}
At inference time the two stages are chained: $\hat{y} = g_y(g_c(x))$. The
two experiments differ substantially in what preprocessing is required
before encoding, and in exactly which data each stage is trained on --
both are described in full below. All code referenced here lives in the
project repository, under \code{emails/scripts/} and \code{cub/scripts/}
respectively; the exact end-to-end evaluation described in
Sections~\ref{sec:emails-eval} and \ref{sec:cub-eval} is implemented in
\code{emails/scripts/report\_test\_accuracy.py} and
\code{cub/scripts/report\_test\_accuracy.py}.

\section{Emails Corpus (Phishing Detection)}

\subsection{Raw data and preprocessing}
\label{sec:emails-preprocessing}

The raw corpus (\code{merged\_emails\_with\_categories.jsonl}, 12{,}300
emails) is \emph{not} used in full -- unlike CUB (Section~\ref{sec:cub-preprocessing}),
the emails pipeline applies a row-level filter before any encoding happens,
because the raw corpus contains material unsuitable for the study
(non-English text, overly long emails, non-human-authored content).

\paragraph{Step 1 -- language and length annotation.} For every email's
\code{Body} text, a word count (whitespace split) and a language guess
(\code{langdetect.detect}, seeded for determinism via
\code{DetectorFactory.seed = 0}) are computed and stored as new columns.
Empty/whitespace-only bodies are assigned \code{word\_count=0},
\code{language="unknown"} without calling the detector.

\paragraph{Step 2 -- filtering to the ``usable'' set.} The three conditions
below are applied simultaneously (logical \textsc{and}, not sequential
steps):

\begin{lstlisting}
usable = original[(original['lang'] == 'en')
                   & (original['words'] < 400)
                   & (original['Created by'] == 'Human')]
\end{lstlisting}

\begin{table}[h]
\centering
\begin{tabular}{ll}
\toprule
Condition & Meaning \\
\midrule
\code{lang == 'en'} & Keep only emails classified as English \\
\code{words < 400} & Keep only emails under 400 words (strict) \\
\code{Created by == 'Human'} & Exclude LLM-generated / synthetic emails \\
\bottomrule
\end{tabular}
\caption{Emails filtering criteria. Applied jointly, 12{,}300 $\rightarrow$ 2{,}330 ``usable'' emails.}
\end{table}

No deduplication is applied to the final data (an abandoned subject-based
rebalancing branch exists earlier in the preprocessing notebook but its
output is not used downstream).

\subsection{Encoding}

Only the \code{Body} text (never \code{Subject}) of each of the 2{,}330
usable emails is encoded, as-is -- no HTML stripping, lowercasing,
whitespace normalization, or additional truncation beyond the word-count
filter above -- using \code{sentence-transformers}' \code{all-MiniLM-L6-v2}
model. Encoding happens once, over the entire usable set, \emph{before} any
train/validation/test split exists.

\subsection{Concept and label construction}

Still on the full, unsplit set: a 19-dimensional binary concept-ground-truth
vector (\code{concept\_gts}) is derived from LLM-labeled emotion/motivation
fields, and reduced to a 6-dimensional filtered vector
(\code{concept\_gts\_f}) matching the six concepts actually used by the
model and shown to human participants in the user study:
\emph{Fear+Authority} (a merge of two raw concepts), \emph{Urgency},
\emph{Curiosity}, \emph{Neutral}, \emph{Reply}, \emph{Open attachment}. The
task label is assigned from the \code{Type} field: \code{Valid}
$\rightarrow 1$, \code{Phishing}/\code{Phishing Simulation} $\rightarrow 0$,
\code{Spam} $\rightarrow 2$ (label $2$ is excluded from CBM training/evaluation,
see Section~\ref{sec:emails-training}).

\subsection{Train / validation / test split}

An earlier attempt at a stochastic, class-balanced split
(\code{random\_state=42}-seeded sampling) was found not to be reproducible
across pandas/numpy versions and is not used. The split actually used is
selected \emph{by fixed email ID} from a pre-recorded list
(\code{archival\_split\_ids.json}), matching the split used in the real
human user study:

\begin{table}[h]
\centering
\begin{tabular}{lrr}
\toprule
Split & Total rows & Rows with label $\in\{0,1\}$ \\
\midrule
Train & 1{,}064 & 551 \\
Validation & 266 & 138 \\
Test (user study set) & 1{,}000 & 1{,}000 \\
\bottomrule
\end{tabular}
\caption{Emails split sizes. The test split is class-balanced by
construction (500 valid / 500 phishing).}
\end{table}

Because the split operates on the already-embedded dataframe, embeddings
are computed exactly once and never recomputed per split.

\subsection{Concept extractor training}
\label{sec:emails-training}

The concept extractor is one binary linear SVM per concept (six
independent classifiers, via scikit-learn's \code{MultiOutputClassifier}),
trained on the \emph{full} training split (all 1{,}064 rows, regardless of
task label -- concept labels are defined for every email, including
\code{Spam}):

\begin{lstlisting}
X_train = np.stack(train['embedding'].values)
y_train = np.stack(train['concept_gts_f'].values)

concept_model = MultiOutputClassifier(
    SVC(kernel='linear', C=1.0)
)
concept_model.fit(X_train, y_train)
\end{lstlisting}

\paragraph{Hyperparameter choice.} A linear kernel with $C=1.0$ was
selected empirically; RBF kernels (both $C=1.0$ and $C=1000$) and a looser
linear SVM ($C=0.01$) were also tried during development and gave lower
validation micro-F1 (documented as code comments alongside the chosen
configuration in the training notebook) -- the reported linear, $C=1.0$
configuration is the one actually used for all downstream results.

\subsection{Label predictor training}

The label predictor is a single multinomial logistic regression, trained
\emph{only} on training rows with a binary task label (\code{train\_class},
551 rows -- \code{Spam} rows are excluded here, since the task being
predicted is binary valid-vs-phishing), using \emph{ground-truth} concepts
rescaled from $\{0,1\}$ to $\{-1,+1\}$:

\begin{lstlisting}
concept_gts_pm1 = 2 * np.stack(train_class['concept_gts_f'].values) - 1
y_train = np.stack(train_class['label'].values)

label_model = LogisticRegression(
    max_iter=1000, C=1, solver='lbfgs',
    fit_intercept=False, penalty=None,
    class_weight='balanced',
)
label_model.fit(concept_gts_pm1, y_train)
\end{lstlisting}

\paragraph{Design choices.} \code{fit\_intercept=False} was chosen so the
fitted weight vector alone is directly interpretable as each concept's
signed contribution to the ``valid email'' decision, with no additive bias
term. \code{penalty=None} (i.e.\ unregularized logistic regression) is
viable here because the input dimensionality is tiny (6 concepts) relative
to the number of training examples (551), so overfitting risk is low;
\code{class\_weight='balanced'} compensates for the training split's class
imbalance (185 valid vs.\ 366 phishing in the training split with binary
labels).

\subsection{End-to-end inference and evaluation}
\label{sec:emails-eval}

At test time, the two trained stages are chained: an email's embedding is
passed through each of the six per-concept SVMs' \emph{decision function}
(continuous distance to the separating hyperplane, not the binarized
prediction), passed through $\tanh$ to bound it to $(-1, 1)$, and the
resulting 6-dimensional vector is fed to the label predictor:

\begin{lstlisting}
logits = np.column_stack(
    [est.decision_function(X_test) for est in concept_model.estimators_]
)
concept_activations = np.tanh(logits)
y_pred = label_model.predict(concept_activations)
\end{lstlisting}

Using the continuous, $\tanh$-transformed decision function (rather than
the SVMs' binarized concept predictions) as the label predictor's input
means the label predictor receives a graded, real-valued confidence signal
per concept at inference time, matching the same $\pm 1$-scaled
representation it was trained on (ground-truth concepts rescaled to
$\{-1,+1\}$ are simply the two extremes of this same range). Evaluated on
the full 1{,}000-row, class-balanced test set, this yields:

\begin{table}[h]
\centering
\begin{tabular}{lrrrr}
\toprule
Class & Precision & Recall & F1 & Support \\
\midrule
Phishing (0) & 0.91 & 0.94 & 0.92 & 500 \\
Valid (1) & 0.93 & 0.91 & 0.92 & 500 \\
\midrule
\textbf{Accuracy} & \multicolumn{4}{r}{\textbf{0.923} (1{,}000 / 1{,}000)} \\
\bottomrule
\end{tabular}
\caption{Emails CBM end-to-end test performance.}
\end{table}

\section{CUB-200-2011 (Bird Species Classification)}

\subsection{Raw data and preprocessing}
\label{sec:cub-preprocessing}

Unlike emails, \textbf{no filtering or exclusion is applied to CUB at
all}: every image in the official CUB-200-2011 train/validation/test split
is encoded, unmodified in selection. Concept and label annotations come
from three pickle files (\code{train.pkl}, \code{val.pkl}, \code{test.pkl},
under \code{class\_attr\_data\_10/}),
the standard 200-class, 312-concept splits used throughout the CBM
literature -- not re-split or re-curated by this study. Loading is a
direct pass-through (CQA's \code{CUBDataset}): the entire pickle file for a
split is loaded in one shot, with no sampling, slicing, deduplication, or
class exclusion.

\subsection{Encoding}

Images are encoded with CLIP (\code{ViT-L/14}), using CLIP's own standard
preprocessing transform, identical for every image and every split:

\begin{lstlisting}
Compose([
    Resize(n_px, interpolation=BICUBIC),
    CenterCrop(n_px),
    convert_to_rgb,
    ToTensor(),
    Normalize(mean=(0.48145466, 0.4578275, 0.40821073),
              std=(0.26862954, 0.26130258, 0.27577711)),
])
\end{lstlisting}

where \code{n\_px} is CLIP's native input resolution (224 for
\code{ViT-L/14}). This is exactly OpenAI's released CLIP preprocessing --
no dataset-specific cropping, augmentation, or normalization is added. The
resulting embeddings are 768-dimensional. Per-split row counts: train
4{,}796, validation 1{,}198, test 5{,}794 (all 200 classes).

\paragraph{Per-split standardization.} Before any model training, each
split's embeddings are independently standardized to zero mean and unit
variance, \emph{per split} (i.e.\ using that split's own mean/std, not
statistics computed on the training split alone):

\begin{lstlisting}
embeddings = (embeddings - embeddings.mean(0, keepdim=True)) \
             / embeddings.std(0, keepdim=True)
\end{lstlisting}

\subsection{Task restriction: the sparrow minimal pair}
\label{sec:cub-restriction}

The CBM experiment does not use the full 200-way classification task.
Instead, both the concept space and the label-prediction task are
restricted to a single, deliberately difficult \emph{minimal pair} of
visually similar species: the Le Conte's Sparrow (class id 123) and the
Savannah Sparrow (class id 126), and to six concepts (of the 112 total)
chosen for their relevance to distinguishing these two species:

\begin{table}[h]
\centering
\begin{tabular}{rl}
\toprule
Concept index & Concept name \\
\midrule
23 & striped breast \\
44 & buff breast \\
48 & white throat \\
69 & buff nape \\
89 & solid belly \\
103 & brown crown \\
\bottomrule
\end{tabular}
\caption{The 6 masked CUB concepts used by the sparrow-pair CBM (indices into the 112-concept CUB concept vector).}
\end{table}

\subsection{Concept extractor training}

\textbf{Crucially, the concept extractor is trained on the training split
with the two sparrow classes excluded} -- it never sees a Le Conte's or
Savannah Sparrow training image, and is evaluated only on the held-out
\emph{test}-split sparrow-pair images (Section~\ref{sec:cub-eval}). This
ensures the concept predictions used at evaluation time come from a model
that generalizes the visual concepts from unrelated species, rather than
one fit directly to the two classes it will later be asked to distinguish:

\begin{lstlisting}
SPARROW_PAIR = [123, 126]
concept_mask = [23, 44, 48, 69, 89, 103]

train_other_subset = np.where(
    ~np.isin(train_y.numpy(), SPARROW_PAIR)
)[0]

concept_model = MultiOutputClassifier(
    SVC(kernel='rbf', C=1.0, class_weight='balanced')
)
concept_model.fit(
    train_embeddings[train_other_subset].numpy(),
    train_concepts[train_other_subset][:, concept_mask].numpy(),
)
\end{lstlisting}

\paragraph{Hyperparameter choice.} An RBF kernel (rather than the linear
kernel used for emails, Section~\ref{sec:emails-training}) was chosen
because CLIP image embeddings and the visual concepts being predicted
(e.g.\ ``striped breast'') are not expected to be linearly separable in
embedding space the way short-text sentence embeddings are for the
email-specific concepts; \code{C=1.0} and \code{class\_weight='balanced'}
match the emails configuration for consistency, compensating for the
imbalance in how often each visual concept is present across the (large,
199-class) training subset used here.

\subsection{Label predictor training}

The label predictor is trained \emph{only} on the sparrow pair's own
training rows (i.e.\ the training images the concept extractor above never
saw), using ground-truth concepts (the same six, masked, rescaled to
$\{-1,+1\}$):

\begin{lstlisting}
train_pair_subset = np.where(
    np.isin(train_y.numpy(), SPARROW_PAIR)
)[0]
train_pair_concepts_pm1 = 2 * train_concepts[train_pair_subset][:, concept_mask].numpy() - 1
train_pair_y = train_y[train_pair_subset].numpy()

label_model = LogisticRegression(
    max_iter=1000, class_weight='balanced'
)
label_model.fit(train_pair_concepts_pm1, train_pair_y)
\end{lstlisting}

Unlike the emails label predictor, no \code{fit\_intercept=False} or
\code{penalty=None} override is used here -- default regularization
(\code{penalty='l2'}, $C=1.0$) and a fitted intercept are used, since with
only two, well-separated classes and 6 ground-truth concepts, the fitted
model reaches perfect training-time separability regardless (see below).

\subsection{End-to-end inference and evaluation}
\label{sec:cub-eval}

Evaluation is restricted to the sparrow pair's rows within the \emph{test}
split (59 images: 29 Le Conte's, 30 Savannah). As with emails, the chain
is: embedding $\rightarrow$ per-concept SVM decision function $\rightarrow$
$\tanh$ $\rightarrow$ label predictor:

\begin{lstlisting}
test_pair_subset = np.where(
    np.isin(test_y.numpy(), SPARROW_PAIR)
)[0]
X_test = test_embeddings[test_pair_subset].numpy()

logits = np.column_stack(
    [est.decision_function(X_test) for est in concept_model.estimators_]
)
concept_activations = np.tanh(logits)
y_pred = label_model.predict(concept_activations)
\end{lstlisting}

\begin{table}[h]
\centering
\begin{tabular}{lrrrr}
\toprule
Class & Precision & Recall & F1 & Support \\
\midrule
Le Conte's Sparrow (123) & 0.78 & 0.86 & 0.82 & 29 \\
Savannah Sparrow (126) & 0.85 & 0.77 & 0.81 & 30 \\
\midrule
\textbf{Accuracy} & \multicolumn{4}{r}{\textbf{0.814} (48 / 59)} \\
\bottomrule
\end{tabular}
\caption{CUB CBM end-to-end test performance on the sparrow minimal pair.}
\end{table}

For reference, using \emph{ground-truth} concepts directly (bypassing the
concept extractor entirely) the label predictor reaches \textbf{100\%}
accuracy on the same 59 test images -- i.e.\ the six chosen concepts are
perfectly sufficient to distinguish the two species when known exactly;
the end-to-end accuracy gap (81.4\% vs.\ 100\%) is attributable entirely
to concept-extraction error on out-of-training-distribution species, not
to the label predictor's decision boundary.

\section{Common implementation details}

Both experiments share the following design choices, applied identically:

\begin{itemize}[leftmargin=*]
    \item \textbf{Concept representation at inference time}: the label
    predictor always consumes $\tanh(\text{decision\_function}(x))$, not
    the SVMs' binarized $\{0,1\}$ predictions -- a continuous, signed
    confidence score in $(-1, 1)$ for each concept.
    \item \textbf{Ground-truth concept scaling at training time}: binary
    ground-truth concepts $\{0,1\}$ are rescaled to $\{-1,+1\}$
    ($2c - 1$) before fitting the label predictor, so that training-time
    inputs (exact $\pm 1$) and inference-time inputs
    ($\tanh$-bounded, approaching but never reaching $\pm 1$) share the
    same representation and sign convention.
    \item \textbf{One-SVM-per-concept}: both concept extractors are a
    scikit-learn \code{MultiOutputClassifier} wrapping an \code{SVC},
    i.e.\ independent binary classifiers per concept (no joint/multi-task
    concept model), fit and queried independently via
    \code{estimator.decision\_function}.
    \item \textbf{Label predictor}: a single multinomial/binomial
    \code{LogisticRegression} in both cases, mapping the (six-dimensional)
    concept-activation vector to the task label.
\end{itemize}

\section{Reproducibility}

\begin{table}[h]
\centering
\begin{tabular}{ll}
\toprule
Artifact & Location \\
\midrule
Emails preprocessing & \code{emails/scripts/preprocessing.ipynb} \\
Emails CBM training + evaluation & \code{emails/scripts/train\_concept\_extractor2.ipynb} \\
Emails accuracy report (standalone) & \code{emails/scripts/report\_test\_accuracy.py} \\
CUB encoding & \code{cub/scripts/encode\_clip.py} \\
CUB CBM training + evaluation & \code{cub/notebooks/train\_model.ipynb} \\
CUB accuracy report (standalone) & \code{cub/scripts/report\_test\_accuracy.py} \\
\bottomrule
\end{tabular}
\caption{Where to find/reproduce each step described in this document.}
\end{table}

Both standalone accuracy-report scripts were verified to reproduce the
exact numbers quoted in this document (0.923 for emails, 0.814 for CUB).

\end{document}
